\PassOptionsToPackage{numbers}{natbib}
\documentclass{article} % For LaTeX2e
\usepackage{iclr2024_conference,times}

\usepackage[utf8]{inputenc} % allow utf-8 input
\usepackage[T1]{fontenc}    % use 8-bit T1 fonts
\usepackage{hyperref}       % hyperlinks
\usepackage{url}            % simple URL typesetting
\usepackage{booktabs}       % professional-quality tables
\usepackage{amsfonts}       % blackboard math symbols
\usepackage{nicefrac}       % compact symbols for 1/2, etc.
\usepackage{microtype}      % microtypography
\usepackage{titletoc}

\usepackage{subcaption}
\usepackage{graphicx}
\usepackage{amsmath}
\usepackage{multirow}
\usepackage{color}
\usepackage{colortbl}
\usepackage{cleveref}
\usepackage{algorithm}
\usepackage{algorithmicx}
\usepackage{algpseudocode}
\usepackage{tikz}
\usepackage{pgfplots}
\usepackage{float}
\usepackage{array}
\usepackage{tabularx}
\pgfplotsset{compat=newest}


\DeclareMathOperator*{\argmin}{arg\,min}
\DeclareMathOperator*{\argmax}{arg\,max}

\graphicspath{{../}} % To reference your generated figures, see below.

\title{Coverage-Gated lDDT-PLI: A Two-Line Change That Decides Whether an Agent Optimising PLINDER Learns Structure or Learns to Abstain}

\author{AIRAS}

\newcommand{\fix}{\marginpar{FIX}}
\newcommand{\new}{\marginpar{NEW}}

\begin{document}

\maketitle

\begin{abstract}
Protein--ligand cofolding benchmarks are increasingly the target of automated research
agents, and on PLINDER the number such an agent will chase is lDDT-PLI. We ask whether
that number is safe to optimise, and show it is not. The reference aggregation in
\texttt{plinder.eval} averages per-ligand lDDT-PLI only over reference ligands that
received an assignable pose, and reports PoseBusters validity as a separate column
rather than gating on it. Both choices are exploitable: abstaining raises the mean, and
a compact object placed anywhere near the pocket centre recovers most reference contacts
because lDDT-PLI is superposition-free by construction. We quantify the damage on 150
PLINDER MLSB test systems using a predictor family whose true quality is fixed by
construction rather than estimated, including an adversary that explicitly searches
orientations to maximise lDDT-PLI. Under the incumbent rule a predictor answering on a
quarter of systems ties the crystal reference at 1.000, and a blob of atoms collapsed
onto a single point outscores a correctly shaped ligand displaced by 4\,\AA{}. We propose
CG-PLI, a two-line change to the aggregation step alone: score unassigned ligands zero,
and gate each per-ligand score on validity relative to the reference pose. Spearman rank
fidelity against the constructed ordering rises from 0.775 to 0.973. All results were
produced on NVIDIA GB200 (Grace-Blackwell, aarch64) and agree bit-for-bit with x86\_64.
\end{abstract}

\section{Introduction}
\label{sec:intro}
Predicting how a small molecule binds a protein is the step that turns structure
prediction into drug discovery, and the last two years have produced both strong
cofolding models \citep{abramson-2024-accurate} and the curated resources needed to evaluate
them. PLINDER \citep{durairaj-2024-plinder} is the most ambitious of these: 449{,}383
protein--ligand systems with similarity-based splits designed to minimise leakage, and
an evaluation package built on OpenStructure that reports lDDT-PLI alongside a
symmetry-corrected RMSD. It is rapidly becoming the default target for method
development in this area.

It is also becoming a target for automation. Autonomous research systems now run the
full loop from literature review through implementation to manuscript
\citep{tang-2025-researcher}, and comparable systems drive physical experiments
\citep{boiko-2023-autonomous}. Those systems are evaluated in settings where the objective is
well posed and difficult to satisfy by accident. A benchmark metric that can be
satisfied by accident behaves very differently once an optimiser is pointed at it, and
this is the situation we examine.

\paragraph{Why this is hard.} The difficulty is not that lDDT-PLI is a bad per-ligand
score. It is a carefully designed one: it compares protein--ligand interatomic distances
rather than absolute coordinates, so it needs no global superposition and a badly placed
ligand cannot corrupt the alignment of everything else. The problem lives one level up,
in how per-ligand scores become the single number that ranks methods. Two properties of
the reference aggregation are individually reasonable and jointly exploitable. First, it
is coverage-blind: \texttt{get\_average\_ligand\_scores} averages over the reference
ligands that received an assignable model pose and skips the rest, so the denominator is
chosen by the predictor. Second, it is validity-blind: PoseBusters
\citep{buttenschoen-2023-posebusters} results are reported beside the score rather than
folded into it, so chemical implausibility and geometric accuracy can be traded against
one another. Neither property is visible in the score itself, and both are invisible to
anything that only reads the leaderboard.

Evidence that this matters already exists but has never been quantified. The MLSB 2024
PLINDER challenge received one genuine submission, at median lDDT-PLI 0.17, and a
deliberately degenerate reference entry reached median 0.67 at a symmetry-corrected RMSD
of 7.93\,\AA{}. That is a single data point on a leaderboard that is no longer
reachable, and it is not enough to tell us how often, or in what direction, the metric
misranks real candidates.

\paragraph{Our contribution.} We turn that anecdote into a measurement and propose a fix.

\begin{itemize}
\item We construct a predictor family whose geometric error and coverage are set by us
      rather than estimated, so that the correct ranking is known a priori and any
      disagreement is attributable to the metric. The family includes an adversary,
      \texttt{blackhole\_busted}, which is a chemically relaxed conformer whose
      orientation is explicitly searched to maximise lDDT-PLI --- a literal model of an
      agent optimising the headline number.
\item We show that on 150 PLINDER MLSB test systems the incumbent aggregation inverts
      the ranking in two distinct, separable ways, and that the inversions are perfectly
      systematic: the bootstrap confidence interval for its rank fidelity has zero width
      across 1{,}000 resamples.
\item We propose CG-PLI, a change confined to the aggregation function that leaves the
      OpenStructure scorer, the PLINDER dataset and every model untouched. It raises
      Spearman rank fidelity from 0.775 to 0.973.
\item We report an incidental result that anyone adopting validity gating needs: only
      77\% of the deposited crystal ligands pass the full PoseBusters suite, so an
      absolute gate penalises the ground truth. Gating relative to the reference pose is
      the necessary correction.
\item We establish that the evaluation stack is numerically portable to NVIDIA
      Grace-Blackwell aarch64, producing bit-identical scores to x86\_64.
\end{itemize}

We verify these claims by running all three aggregation rules over the same per-system
scores on one GB200 node, with 1{,}000-replicate bootstrap intervals over systems.

\section{Related Work}
\label{sec:related}
\paragraph{Benchmarks for protein--ligand pose prediction.} PLINDER
\citep{durairaj-2024-plinder} is the resource we modify. Its contribution is a
similarity-clustered split that de-leaks PLINDER train from PLINDER test, plus an
evaluation package wrapping OpenStructure's lDDT-PLI and symmetry-corrected RMSD
scorers. Our work is not an alternative benchmark and does not dispute the split; we
take the dataset and the per-ligand scorer as given and change only the function that
reduces per-ligand scores to a per-method number. Runs N' Poses
\citep{skrinjar-2025-have} addresses a different weakness of the same family of evaluations,
namely that most of PLINDER test predates the training cutoff of current cofolding
models, and supplies a post-cutoff benchmark with published per-method lDDT-PLI and
coverage. Notably, the coverage figures reported there range from 71\% to 96\% across
methods, which is precisely the regime in which coverage-blind averaging distorts a
comparison; that work reports coverage alongside the mean but does not fold it in, which
is the step we take. PoseBench \citep{morehead-2025-assessing} demonstrates that method
rankings change depending on which evaluation axis is chosen, and argues for multi-axis
reporting. We agree with the diagnosis and differ in the remedy: reporting more columns
leaves the reader to combine them, whereas we argue that coverage and validity are not
optional columns but corrections that belong inside the headline statistic.

\paragraph{Physical validity as an evaluation axis.} PoseBusters
\citep{buttenschoen-2023-posebusters} established that geometric accuracy alone cannot
certify a docking method, since deep learning methods routinely emit poses with
impossible bond lengths, non-planar aromatic rings and steric clashes. It is the source
of our validity gate, and its architecture-neutral pure-Python implementation is what
makes the gate portable. Interaction-fingerprint analysis
\citep{errington-2025-assessing} adds a third axis, showing that models with acceptable
geometric scores often fail to recover the specific interactions that matter. Our
position is a strict extension of both: we use validity not as a reported column but as
a per-ligand multiplier, so that it cannot be traded off against contact recovery. We do
not incorporate interaction fingerprints, and doing so is a natural extension.

\paragraph{Adversarial probing of structure predictors.} Prior work has stress-tested
cofolding models with physically and chemically motivated perturbations and found
divergence from expected physical behaviour that aggregate accuracy hides. That work
perturbs the \emph{model} to expose what the model has learned. We invert the target:
our perturbations construct \emph{predictors} of known quality in order to expose what
the \emph{metric} has learned to reward. To our knowledge no prior work measures rank
fidelity of a protein--ligand evaluation metric against a constructed ground-truth
ordering.

\paragraph{Autonomous research agents.} AI-Researcher \citep{tang-2025-researcher} and
Coscientist \citep{boiko-2023-autonomous} demonstrate end-to-end autonomous scientific work.
Both are evaluated on objectives that are hard to game --- code that must run and pass
tests, or a chemical yield that must be physically realised. Neither addresses what
happens when the objective admits a degenerate optimum. Our result is a concrete
instance of that gap in a domain those systems are actively being pointed at, and the
mitigation we propose is cheap enough to adopt before it becomes a problem.

\section{Background}
\label{sec:background}
\paragraph{Problem setting.} Fix an evaluation set of systems indexed by $j$. System $j$
carries a set of reference ligands $R_j$; in the subset we use, $|R_j| = 1$ throughout.
A predictor is a map from a system to a set of candidate poses, and it may return
nothing for a given reference ligand. Scoring proceeds in two stages that this paper
carefully separates.

The first stage is per-ligand and is not our subject. For reference ligand $i$, an
assignment step matches a model pose to it by graph isomorphism with symmetry
correction, and a scorer produces $s_i \in [0,1]$. When no model pose can be assigned,
$s_i$ is undefined and we write $s_i = \emptyset$. The scorer is lDDT-PLI. Let
$C_i = \{(p,\ell) : d^{\text{ref}}(p,\ell) < r\}$ be the set of receptor-atom /
ligand-atom pairs within an inclusion radius $r = 6.0$\,\AA{} in the reference complex.
Then
$$
s_i \;=\; \frac{1}{|T|}\sum_{\tau \in T}
\frac{\bigl|\{(p,\ell) \in C_i : |\,d^{\text{mdl}}(p,\ell) - d^{\text{ref}}(p,\ell)\,| < \tau\}\bigr|}{|C_i|},
\qquad T = \{0.5, 1, 2, 4\}\,\text{\AA{}} .
$$
The two features of this definition that drive everything below are that it depends only
on \emph{differences of interatomic distances}, so no superposition is required, and
that it is an average over tolerances of which half are coarse. A pose need not be in
the right place to preserve a distance to within 4\,\AA{}; it needs only to be at
roughly the right remove from the same receptor atoms.

The second stage is aggregation, and it is the object of study. An aggregation rule $A$
maps the vector $(s_i)_{i \in R}$ over all reference ligands in the evaluation set, plus
whatever side information it chooses to use, to a scalar. The incumbent rule implemented
by \texttt{get\_average\_ligand\_scores} is
$$
\mathrm{PLI}_{\text{ave}} \;=\; \frac{1}{|\{i : s_i \neq \emptyset\}|}\sum_{i \,:\, s_i \neq \emptyset} s_i ,
$$
with a heavy-atom-weighted variant $\mathrm{PLI}_{\text{wave}}$ that weights each term by
the ligand's heavy-atom count. Both condition on the set of ligands the predictor chose
to answer.

\paragraph{Companion axes.} Two further quantities are standard. The symmetry-corrected
RMSD between model and reference ligand heavy atoms measures placement directly and is
sensitive exactly where lDDT-PLI saturates; because the receptor is identical between
model and reference in our setting, no superposition is applied and this reduces to an
in-place symmetry-minimised RMSD. PoseBusters supplies a binary validity verdict
$v_i \in \{0,1\}$ from a conjunction of chemical and geometric checks
\citep{buttenschoen-2023-posebusters}. Coverage, the fraction of reference ligands that
received an assignable pose, is what $\mathrm{PLI}_{\text{ave}}$ silently conditions on.

\paragraph{An unusual assumption we make explicit.} To ask whether a metric ranks
correctly, one needs a correct ranking to compare against. We obtain it by construction
rather than by measurement, defining the utility of a predictor as
$U = \text{coverage} \times \exp(-\mathrm{rmsd}/2) \times \text{validity}$. This encodes
a normative claim --- a method that answers half the time is worth half as much, and a
chemically impossible pose is worth nothing --- which we state here rather than derive,
because the reader should be able to reject it independently of our measurements.

\section{Method}
\label{sec:method}
CG-PLI changes the aggregation step and nothing else. The per-ligand scorer, the
assignment procedure, the inclusion radius, the tolerance set and the dataset are all
left exactly as PLINDER defines them. Written against the notation above, the proposed
statistic is
$$
\mathrm{CG\text{-}PLI} \;=\; \frac{1}{|R|}\sum_{i \in R} v_i \, \tilde{s}_i ,
\qquad
\tilde{s}_i = \begin{cases} s_i & s_i \neq \emptyset \\ 0 & s_i = \emptyset \end{cases}
$$
where the sum now runs over \emph{every} reference ligand in the evaluation set and
$v_i$ is the PoseBusters indicator. The two modifications are independent and address
disjoint failure modes.

\emph{Coverage gating} replaces the predictor-chosen denominator with $|R|$ and scores
unanswered ligands zero. The motivation is that lDDT-PLI is a recall-style quantity
evaluated on a support the predictor selects, and any recall-style score restricted to
the cases where a system chose to answer is maximised by answering less. This is exactly
the pathology that the selective-prediction literature addresses, and the standard
remedy --- convert a conditional mean into an unconditional one --- transfers without
modification. Note that this does not merely penalise abstention, it removes the
incentive entirely: under CG-PLI a predictor can never improve its score by withholding
a pose, because a withheld pose contributes zero and an emitted one contributes at worst
zero.

\emph{Validity gating} multiplies each per-ligand score by $v_i$, so a chemically
implausible pose contributes nothing regardless of how many reference contacts it
happens to reproduce. Placing the gate at the level of the individual ligand rather than
reporting an aggregate validity rate is what prevents the two axes being traded against
each other: a method cannot compensate for invalid poses on some systems by recovering
more contacts on others.

The composite has a property worth stating because it makes adoption safe. Since
$v_i \in \{0,1\}$, $\tilde{s}_i \le s_i$ and $|R| \ge |\{i : s_i \neq \emptyset\}|$,
we have $\mathrm{CG\text{-}PLI} \le \mathrm{PLI}_{\text{ave}}$ always, with equality
exactly when coverage is complete and every emitted pose is valid. CG-PLI is therefore a
lower bound on the incumbent score that coincides with it precisely for the methods that
have nothing to hide, and reporting both costs one extra column.

\paragraph{Gating relative to the reference.} A naive implementation uses the absolute
PoseBusters verdict, and we found that this is wrong in a way that matters. Deposited
crystal ligands are strained relative to a generated conformer ensemble and frequently
fail the energy-ratio check, among others; in our evaluation set only 77\% of reference
ligands pass the full suite. An absolute gate therefore scores the ground truth at 0.77
rather than 1.00, which is indefensible for a metric whose upper anchor is the
experimental structure. We therefore gate relative to the reference: writing $c_k(i)$
for the outcome of check $k$ on pose $i$ and $c_k(\mathrm{ref})$ for the same check on
the reference pose of that system,
$$
v_i \;=\; \prod_{k \,:\, c_k(\mathrm{ref}) = 1} c_k(i) ,
$$
so a pose is penalised only for failing a check that the reference itself passes. This
restores the reference to exactly 1.000 by construction while retaining full
discriminative power against the degenerate predictors, which fail checks the reference
passes.

\begin{figure}[t]
\centering
\includegraphics[width=0.85\linewidth]{images/diagram/cgpli_method.pdf}
\caption{The CG-PLI aggregation. The per-ligand OpenStructure scorer is untouched; the two gates act only where per-ligand scores are combined. Higher aggregate values indicate better predictors.}
\label{fig:method}
\end{figure}

\section{Experimental Setup}
\label{sec:experimental}
\paragraph{Data.} We use the MLSB subset of the PLINDER 2024-06/v2 test split
\citep{durairaj-2024-plinder}, taking the first 150 systems in shard order from the 346
single-ligand systems listed in \texttt{splits/mlsb/inputs.parquet}. Every system is
fetched anonymously over HTTPS from the public bucket; no cloud credentials or project
are required. From each system we take the receptor heavy atoms parsed from
\texttt{receptor.pdb} and the reference ligand from the single SDF in
\texttt{ligand\_files/}. Reference contacts are computed once per system at a 6.0\,\AA{}
inclusion radius and reused across every predictor, which is what makes the per-predictor
scores directly comparable.

\paragraph{The predictor family.} Because the question is whether a metric ranks
correctly, we need candidates whose ranking is known rather than estimated. We therefore
construct eleven predictors whose geometric error and coverage are set by us.
\emph{reference} emits the holo ligand and anchors the top of the scale.
\emph{perturbed-$k$} for $k \in \{0.5, 1, 2, 4, 8\}$\,\AA{} applies a rigid translation
of magnitude $k$ in a random direction; because the translation is rigid, the
symmetry-corrected RMSD is exactly $k$, so this family is a calibration ladder with no
estimation error. \emph{decoy} generates an independent ETKDG conformer of the correct
molecule, relaxes it with MMFF, and places it at the pocket centroid in one arbitrary
orientation; it is chemically valid but carries no information about the true pose.
\emph{blackhole} collapses every heavy atom onto the pocket centroid, retaining a
jitter of 0.05\,\AA{} only to avoid exactly coincident atoms. \emph{blackhole\_busted}
is the adversary: it takes the same relaxed conformer as \emph{decoy} and searches 16
random orientations at the pocket centroid, keeping the one that maximises lDDT-PLI. It
is a direct model of an agent that has been told to maximise the headline number and has
a small compute budget with which to do so. Finally \emph{abstain-$p$} for
$p \in \{0.25, 0.5\}$ emits the reference pose on the easiest fraction $p$ of systems,
ranked by ligand heavy-atom count, and nothing elsewhere; its poses are perfect where it
answers, so it isolates coverage gaming from pose quality. The pocket centroid used by
the degenerate predictors is computed from receptor atoms lining the reference ligand,
which is information a pocket-conditioned predictor legitimately has.

\paragraph{Metrics and comparison.} The primary metric is the Spearman rank correlation
between the ordering an aggregation rule induces over the eleven predictors and the a
priori utility ordering defined in the Background. We compare three rules over identical
per-system scores: the proposed CG-PLI and the two incumbents
$\mathrm{PLI}_{\text{ave}}$ and $\mathrm{PLI}_{\text{wave}}$. Including the weighted
variant is an ablation on the aggregation: it shares the coverage blindness and the
validity blindness while differing in weighting, so if the failure were caused by the
unweighted mean rather than by the conditioning, the two incumbents would separate.
Alongside the primary metric we report per-predictor lDDT-PLI, symmetry-corrected RMSD,
PoseBusters pass rate both absolute and relative to the reference, and coverage.
Uncertainty comes from a 1{,}000-replicate bootstrap that resamples systems with
replacement and recomputes every predictor's aggregate and the resulting rank
correlation on each replicate.

\paragraph{Implementation and hardware.} Each of the three aggregation rules runs as a
separate job that recomputes the per-system scores from scratch rather than reading a
shared cache, so no two runs can report the same numbers for the wrong reason. Runs were
executed on a single NVIDIA GB200 (Grace-Blackwell, aarch64) node at RIKEN R-CCS through
a Slurm backend, inside an Apptainer container built from a pinned dependency set with a
committed lock file. Scoring is parallelised across systems with 30 worker processes and
PoseBusters is invoked once per system for all eleven poses; 150 systems complete in
61.5\,s and a whole run finishes in under nine minutes. The workload is CPU-bound
cheminformatics and requires no CUDA. Random seeds are fixed per system as a base seed
plus the system index, so the perturbation directions and orientation searches are
reproducible.

\section{Results}
\label{sec:results}
\subsection{Rank fidelity}

Over the 150 evaluated systems, CG-PLI attains a Spearman rank correlation of
\textbf{0.9725} against the constructed utility ordering, with a 95\% bootstrap interval
of $[0.9083, 0.9815]$. Both incumbent rules attain \textbf{0.7747}. The gap in favour of
the proposed rule is $+0.1978$ and the intervals are disjoint. \Cref{fig:spearman} shows the
three rules with their bootstrap intervals.

The interval on the incumbent deserves comment: it is $[0.7747, 0.7747]$, of zero width
across 1{,}000 resamples. The ordering that $\mathrm{PLI}_{\text{ave}}$ induces over the
predictor family did not change under any resample of the systems. Its failure is
therefore systematic rather than a small-sample artefact, and collecting more PLINDER
systems would not repair it.

$\mathrm{PLI}_{\text{wave}}$ produces different values from $\mathrm{PLI}_{\text{ave}}$
--- 0.481 against 0.502 for \emph{blackhole}, for instance --- but an identical ordering
and therefore an identical rank correlation. This ablation isolates the cause: the
failure is the conditioning on the answered subset together with the missing validity
gate, not the choice of weighting.

\subsection{Where the incumbent inverts}

Table~\ref{tab:main} gives the full per-predictor breakdown, and \Cref{fig:bypredictor} plots
the same data. Two inversions are visible and they are separable.

\begin{table}[t]
\centering
\small
\caption{Per-predictor results on 150 PLINDER MLSB systems. For utility, CG-PLI,
$\mathrm{PLI}_{\text{ave}}$, $\mathrm{PLI}_{\text{wave}}$, coverage and both PoseBusters
columns, higher is better; for BiSyRMSD, lower is better. Rows are ordered by the a
priori utility. Bold marks the two inversions discussed in the text.}
\label{tab:main}
\begin{tabular}{lrrrrrrrr}
\toprule
Predictor & Utility & CG-PLI & PLI$_{\text{ave}}$ & PLI$_{\text{wave}}$ & Cov. & PB abs. & PB rel. & RMSD (\AA{}) \\
\midrule
reference          & 1.000 & 1.000 & 1.000 & 1.000 & 1.00 & 0.77 & 1.00 & 0.00 \\
perturbed-0.5      & 0.779 & 0.740 & 1.000 & 1.000 & 1.00 & 0.53 & 0.74 & 0.50 \\
perturbed-1        & 0.607 & 0.281 & 0.873 & 0.873 & 1.00 & 0.19 & 0.32 & 1.00 \\
abstain-0.5        & 0.500 & 0.500 & \textbf{1.000} & 1.000 & 0.50 & 0.76 & 1.00 & 0.00 \\
perturbed-2        & 0.368 & 0.041 & 0.689 & 0.689 & 1.00 & 0.03 & 0.06 & 2.00 \\
abstain-0.25       & 0.250 & 0.253 & \textbf{1.000} & 1.000 & 0.25 & 0.71 & 1.00 & 0.00 \\
perturbed-4        & 0.135 & 0.018 & 0.463 & 0.462 & 1.00 & 0.02 & 0.05 & 4.00 \\
perturbed-8        & 0.018 & 0.012 & 0.197 & 0.198 & 1.00 & 0.09 & 0.12 & 8.00 \\
decoy              & 0.000 & 0.004 & 0.395 & 0.368 & 1.00 & 0.00 & 0.01 & 6.03 \\
blackhole\_busted  & 0.000 & 0.005 & \textbf{0.555} & 0.532 & 1.00 & 0.00 & 0.01 & 4.33 \\
blackhole          & 0.000 & 0.000 & \textbf{0.502} & 0.481 & 1.00 & 0.00 & 0.00 & 4.92 \\
\bottomrule
\end{tabular}
\end{table}

\paragraph{Coverage gaming.} \emph{abstain-0.25} answers on a quarter of the systems and
is silent on the rest, yet $\mathrm{PLI}_{\text{ave}}$ scores it 1.000, exactly level
with the crystal reference. \emph{abstain-0.5} likewise scores 1.000. Under the incumbent
rule, abstention is not merely unpunished: it is the cheapest available route to a
perfect score, since a predictor that answers only where it is confident inherits the
score of its best cases. CG-PLI assigns 0.253 and 0.500, matching the constructed
utilities of 0.25 and 0.50 to within 0.003.

\paragraph{Degenerate placement.} \emph{blackhole} --- every heavy atom collapsed onto a
single point, median RMSD 4.92\,\AA{}, zero PoseBusters validity --- scores 0.502 under
$\mathrm{PLI}_{\text{ave}}$, above \emph{perturbed-4} at 0.463, which is a correctly
shaped ligand displaced by exactly 4\,\AA{}. The adversary \emph{blackhole\_busted}
reaches 0.555, outranking both \emph{perturbed-4} and \emph{perturbed-8}, and
\emph{decoy} at 0.395 also outranks \emph{perturbed-8}. Three predictors that carry no
information about the true pose beyond the location of the pocket thus place above a
predictor with bounded 4\,\AA{} error. This is the MLSB 2024 observation reproduced
systematically rather than as a single leaderboard row, and the adversary shows the
effect is reachable by search: 16 random orientations were enough to lift a meaningless
conformer from 0.395 to 0.555.

The mechanism is the one anticipated in the Background. Because lDDT-PLI scores
differences of interatomic distances and averages over tolerances of which half are
2\,\AA{} and 4\,\AA{}, a compact object near the pocket centre preserves a large
fraction of reference contact distances. Contact recovery saturates well before
placement is correct.

\subsection{Which gate does the work, and where the method falls short}

The two gates are not redundant. The abstainers emit perfectly valid poses, with relative
PoseBusters rates of 1.00, so validity gating leaves them untouched and only coverage
gating corrects them. The degenerate predictors have full coverage, so coverage gating
leaves them untouched and only validity gating corrects them, driving all three to at
most 0.005 given relative pass rates of 0.00--0.01. Removing either gate would leave one
inversion intact.

\Cref{fig:utility} plots each rule against the constructed utility and makes the residual
weakness of CG-PLI visible. The proposed rule is monotone but markedly steeper than
utility in the mid range: \emph{perturbed-1} has utility 0.607 against CG-PLI 0.281, and
\emph{perturbed-2} has utility 0.368 against 0.041. The cause is the validity gate rather
than coverage. Relative PoseBusters pass rates fall from 0.74 at 0.5\,\AA{} displacement
to 0.32 at 1\,\AA{} and 0.06 at 2\,\AA{}, because rigidly translating a crystal ligand
within its own pocket induces protein--ligand clashes almost immediately. PoseBusters is
right to reject those poses, but the resulting aggregate understates the residual
usefulness of a near miss. Since the primary metric is rank-based this does not affect
our headline claim, but anyone using CG-PLI as a continuous score should understand that
it behaves closer to a fraction-of-systems-solved-validly than to an average contact
recovery. Softening the hard conjunction into a weighting by the number of failed checks
is the obvious next step.

A second limitation is that the validity gate inherits PoseBusters' own calibration. The
77\% absolute pass rate on deposited crystal ligands is the clearest evidence that this
calibration is not perfect, and it is why the relative gate is necessary rather than
merely convenient. A third is that our headline comparison depends on the constructed
utility function, whose form is a normative choice we stated in advance rather than
derived; a reader who rejects the claim that coverage should multiply linearly would
draw a different, though still monotone, conclusion.

\begin{figure}[t]
\centering
\includegraphics[width=0.85\linewidth]{images/comparison/comparison_spearman_correlation.pdf}
\caption{Spearman rank fidelity of each aggregation rule against the constructed utility ordering, with 95\% bootstrap intervals over 1{,}000 resamples of the 150 systems. Higher is better. The interval on both incumbent rules has zero width.}
\label{fig:spearman}
\end{figure}

\begin{figure}[t]
\centering
\includegraphics[width=0.85\linewidth]{images/comparison/comparison_score_by_predictor.pdf}
\caption{Aggregate score assigned to each predictor by each rule, with the a priori utility overlaid. Higher is better for every series. Crossings between a rule's curve and the utility curve are ranking inversions.}
\label{fig:bypredictor}
\end{figure}

\begin{figure}[t]
\centering
\includegraphics[width=0.85\linewidth]{images/chart/aggregate_vs_utility.pdf}
\caption{Aggregate score against a priori utility for the proposed and incumbent rules. Higher is better on both axes; a faithful metric is monotone in utility, and the dashed line is the identity.}
\label{fig:utility}
\end{figure}

\subsection{Cross-architecture reproducibility}

The sanity configuration produced a rank correlation of 0.6482912794051882 on both
x86\_64 and GB200 aarch64, identical to all reported digits. We are not aware of prior
evidence that this evaluation stack is numerically portable to Grace-Blackwell, and it is
a precondition for running PLINDER-based evaluation on such systems at all. The claim is
limited: it covers our own scorer implementation on the synthetic sanity systems, not
OpenStructure itself, whose aarch64 behaviour remains to be verified.

\section{Conclusion}
\label{sec:conclusion}
We asked whether lDDT-PLI, as aggregated by PLINDER's evaluation package, is a safe
objective for an automated research agent, and found that it is not. On 150 PLINDER MLSB
test systems the incumbent aggregation assigns a perfect score to a predictor that
answers on a quarter of the systems, and ranks a blob of atoms collapsed onto a single
point above a correctly shaped ligand displaced by 4\,\AA{}. Both failures are properties
of the aggregation step rather than of the per-ligand score, and both are perfectly
systematic: the bootstrap interval on the incumbent's rank fidelity has zero width.

The proposed correction, CG-PLI, is deliberately small. Score unassigned reference
ligands zero instead of skipping them, and multiply each per-ligand score by a validity
indicator taken relative to the reference pose. Nothing else changes --- not the
OpenStructure scorer, not the assignment procedure, not the dataset, and no model is
retrained. Rank fidelity rises from 0.775 to 0.973, the two gates provably address
disjoint failure modes, and because CG-PLI never exceeds the incumbent and coincides with
it exactly when coverage is complete and all poses are valid, it can be reported
alongside the existing number at the cost of one column.

Two findings along the way are worth carrying forward independently of the proposal.
Only 77\% of deposited crystal ligands pass the full PoseBusters suite, so any validity
gating scheme must be defined relative to the reference or it will penalise the ground
truth. And the evaluation stack reproduces bit-identically between x86\_64 and NVIDIA
Grace-Blackwell aarch64, which removes an unverified assumption for anyone intending to
run these benchmarks on such hardware.

The obvious extensions are to soften the hard validity conjunction into a graded penalty,
which would remove the excessive steepness we observe in the mid range; to fold
interaction-fingerprint recovery in as a third gate \citep{errington-2025-assessing}; and to
apply the rank-fidelity protocol itself to other benchmarks. That last direction is the
one we would prioritise. The methodology here --- construct predictors of known quality,
including an adversary that explicitly optimises the metric, then measure how faithfully
the metric ranks them --- is not specific to protein--ligand docking. As autonomous
agents \citep{tang-2025-researcher,boiko-2023-autonomous} are pointed at more benchmarks,
checking that a headline number cannot be satisfied by accident should become a routine
step in publishing one.

This work was generated by \textsc{AIRAS} \citep{airas2025}.

\bibliographystyle{iclr2024_conference}
\bibliography{references}

\end{document}
